\documentclass[11pt]{amsart}
\usepackage{amsmath,amssymb,amsthm}
\usepackage[margin=1.15in]{geometry}
\usepackage[hidelinks]{hyperref}

\newtheorem{theorem}{Theorem}[section]
\newtheorem{lemma}[theorem]{Lemma}
\theoremstyle{remark}
\newtheorem{remark}[theorem]{Remark}

\newcommand{\doi}[1]{\href{https://doi.org/#1}{doi:\nolinkurl{#1}}}
\newcommand{\arxiv}[1]{\href{https://arxiv.org/abs/#1}{arXiv:\nolinkurl{#1}}}
\newcommand{\pd}{\partial}
\setlength{\emergencystretch}{2.5em}

\begin{document}

\title[Connective constant of non-simple transitive graphs of degree five]{The connective constant of
non-simple transitive graphs of degree five}
\author{Yitzchak Shmalo}
\address{Einstein Institute of Mathematics, The Hebrew University of Jerusalem, Jerusalem, Israel}
\email{yspennstate@gmail.com}
\subjclass[2020]{05C30, 05C40, 60K35, 82B41}
\keywords{self-avoiding walk, connective constant, transitive graph, multigraph, edge connectivity}

\begin{abstract}
Grimmett and Li proved that the connective constant $\mu$ of an infinite, connected, $\Delta$-regular,
vertex-transitive graph satisfies $\mu\ge\sqrt{\Delta-1}$ when the graph is simple, and also when it has multiple
edges and $\Delta\le4$, and asked whether the second statement holds for $\Delta\ge5$. We prove it for $\Delta=5$:
every infinite, connected, $5$-regular, vertex-transitive, loopless multigraph has $\mu\ge2$, with equality for the
bridge graph. The proof passes to the quotient by the atoms, the smallest finite vertex sets with the fewest boundary edges, where the branching
inequality of Grimmett and Li holds. Each walk of the quotient lifts to many walks of the graph, counted by
self-avoiding paths inside the atoms, and a convexity argument gives growth rate at least $\sqrt{\Delta-1}$ under
a condition on the atoms alone. For degree five the four possible atom types satisfy it.
\end{abstract}

\maketitle

\section{Introduction}\label{sec:intro}

Let $G=(V,E)$ be an infinite, connected, locally finite graph, possibly with multiple edges but without loops. A
self-avoiding walk (SAW) is an alternating sequence of vertices and edges in which consecutive vertices are the ends
of the edge between them and no vertex appears twice; parallel edges give different walks. Fix $v\in V$ and let
$c_k(v)$ be the number of $k$-step SAWs starting at $v$. When $G$ is vertex-transitive, $c_k=c_k(v)$ does not depend
on $v$, the sequence is submultiplicative, and $\mu(G)=\lim_kc_k^{1/k}=\inf_kc_k^{1/k}$ is the \emph{connective
constant} of $G$ \cite{HM54}. If $G$ is $\Delta$-regular, that is, every vertex has degree $\Delta$, then $\mu(G)\le\Delta-1$.

Grimmett and Li proved the following lower bound \cite[Theorem~4.1]{GL15}. Let $G$ be an infinite, connected,
$\Delta$-regular, vertex-transitive graph with $\Delta\ge2$. If $G$ is simple, then $\mu(G)\ge\sqrt{\Delta-1}$;
if $G$ is non-simple and $\Delta\le4$, then again $\mu(G)\ge\sqrt{\Delta-1}$. The bound is attained by the bridge
graph $B_\Delta$, the two-way infinite path in which every second edge is replaced by $\Delta-1$ parallel edges. They
ask whether the second statement extends to $\Delta\ge5$ \cite{GL15}; this is Question~2.2 of their survey
\cite{GL19}, and Question~4 of \cite{GL13}. We settle the first open case.

\begin{theorem}\label{thm:main}
Let $G$ be an infinite, connected, $5$-regular, vertex-transitive, loopless multigraph. Then $\mu(G)\ge2$.
\end{theorem}

The bridge graph $B_5$ has $\mu(B_5)=2$, so the constant cannot be improved.

\subsection*{The method}
Let $\lambda(G)$, the \emph{finite-set edge connectivity}, be the least number of edges leaving a finite nonempty set
of vertices. If $\lambda(G)=\Delta$,
the inequality follows from the branching argument of Grimmett and Li, which needs $\Delta$ edge-disjoint infinite
SAWs from every vertex (Section~\ref{sec:branch}). The difficulty with multiple edges is that $\lambda(G)$ may be
smaller than $\Delta$; in $B_\Delta$ it is $2$. When $\lambda<\Delta$, the finite sets of least size attaining
$\lambda$, the \emph{atoms}, partition $V$ into blocks of a common size $a$, and every vertex sends the same number $e$
of edges out of its block (Section~\ref{sec:atoms}; these are edge versions of the atoms used by Watkins \cite{Wat70} for the
connectivity of finite transitive graphs). Mader \cite{Mad71} proved that the edge connectivity of a finite connected
vertex-transitive simple graph equals its degree. Collapsing each atom to a vertex gives an $ae$-regular quotient
$H$ in which every finite set has at least $ae$ boundary edges, so the branching inequality holds in $H$. A SAW of $H$
lifts to SAWs of $G$ by choosing a SAW inside each atom it visits, and the lifts of a fixed walk are counted by a
product of path sums (Section~\ref{sec:lift}). Lemma~\ref{lem:weighted} gives a criterion, in terms of these path sums
alone, under which the branching inequality forces the weighted count of lifts to stay bounded below at the scale
$(\Delta-1)^{-k/2}$, which gives $\mu\ge\sqrt{\Delta-1}$. The criterion reads
\[
\psi\bigl(\tfrac{ae-2}2\bigr)\ \ge\ 0,
\]
where $W(b)$ is the least total weight of $b+1$ of the edges by which a lifted walk can leave an atom, and $\psi$ is
the lower convex envelope of $\log W$ on $\{0,1,\dots,ae-2\}$. For $\Delta=5$ there are four atom types, and
Section~\ref{sec:five} checks the criterion for each.

\begin{center}
\begin{tabular}{ll}
Section~\ref{sec:branch} & edge-disjoint rays and the branching inequality (Lemma~\ref{lem:rays}, Lemma~\ref{lem:branch})\\
Section~\ref{sec:atoms} & atoms and the quotient $H$ (Lemma~\ref{lem:atoms})\\
Section~\ref{sec:lift} & lifting walks of $H$ to $G$; the weighted criterion (Lemma~\ref{lem:weighted})\\
Section~\ref{sec:five} & the four atom types of degree five; proof of Theorem~\ref{thm:main}\\
Section~\ref{sec:remarks} & what the criterion gives in higher degree
\end{tabular}
\end{center}

Throughout, $T=\sqrt{\Delta-1}$ and $t=1/T$. Since $\mu=\inf_kc_k^{1/k}$, the power series $\sum_kc_kz^k$ has radius
of convergence $1/\mu$; so $\sum_kc_kt^k=\infty$ implies $\mu\ge T$.

\section{Edge-disjoint rays and the branching inequality}\label{sec:branch}

In this section $H$ is an infinite, connected, loopless, $d$-regular multigraph with $d\ge2$ such that every finite
nonempty vertex set $S$ satisfies $|\pd S|\ge d$, where $\pd S$ is the set of edges with exactly one end in $S$. A
\emph{ray} is an infinite SAW.

\begin{lemma}\label{lem:rays}
For every vertex $w$ of $H$ there are $d$ edge-disjoint rays from $w$, one beginning with each edge at $w$.
\end{lemma}

\begin{proof}
Let $R$ be large and let $H_R$ be the finite multigraph obtained from $H$ by identifying all vertices at graph
distance more than $R$ from $w$ to a single vertex $z$ and deleting the loops this creates. A set of edges separating
$w$ from $z$ in $H_R$ contains $\pd S$ for a finite set $S\ni w$ of vertices of $H$, so it has at least $d$ edges. By
Menger's theorem there are $d$ edge-disjoint paths from $w$ to $z$ in $H_R$; truncating each before $z$ gives a path in $H$
from $w$ reaching distance $R$. Since $H$ is locally finite, a diagonal argument over $R\to\infty$ gives $d$ edge-disjoint rays from
$w$. As $w$ has exactly $d$ edges, each ray begins with a different one.
\end{proof}

Let $\pi=v_0h_0v_1h_1v_2\cdots$ be a ray of $H$, and write $\pi_s$ for its initial segment ending at $v_s$. For
$s\ge0$, an edge $f$ at $v_s$, not on $\pi_s$, is \emph{extendable at $s$} if some ray begins with $\pi_s$ followed
by $f$; the edge $h_s$ is one of them. Write $B_s$ for the number of edges other than $h_s$ that are extendable at
$s$. This number depends only on $\pi_s$. By Lemma~\ref{lem:rays} every edge at $v_0$ is extendable at $0$, so
$B_0=d-1$. For $s\ge1$ there are $d-2$ edges at $v_s$ other than $h_{s-1}$ and $h_s$, and we write $b_s=B_s$ and
$r_s=d-2-b_s$ for the numbers of these that are and are not extendable at $s$. Put $b_0=d-2$.

The following is inequality (5.3) of \cite{GL15}, which Grimmett and Li derive from their condition $\Pi$, a rule assigning to each
non-extendable edge a path back to the walk that begins with that edge. We give
the argument directly in terms of edge-disjoint rays; it includes the case of an edge that leads straight back to an
earlier vertex of $\pi$, which condition $\Pi$, read literally, does not cover.

\begin{lemma}\label{lem:branch}
For every ray $\pi$ of $H$ and every $n\ge1$,
\[
\sum_{s=0}^{2n-1}b_s\ \ge\ n(d-2).
\]
\end{lemma}

\begin{proof}
Let $N=2n$, $B=\sum_{s=0}^{N-1}B_s$ and $R=\sum_{s=1}^{N-1}r_s$. Then
\begin{equation}\label{eq:BR}
B+R=(d-1)+(N-1)(d-2)=N(d-2)+1.
\end{equation}
We construct an injective map from the non-extendable edges counted by $R$ to the extendable edges counted by $B$.

Let $1\le s\le N-1$. By Lemma~\ref{lem:rays} choose $d$ edge-disjoint rays from $v_s$, one beginning with each edge at
$v_s$. Let $g$ be an edge at $v_s$, other than $h_{s-1}$ and $h_s$, that is not extendable at $s$, and let $\rho$ be
the chosen ray beginning with $g$. Then $\rho$ meets $\{v_0,\dots,v_{s-1}\}$, since otherwise $\pi_s$ followed by
$\rho$ would be a ray. Let $y$ be the first vertex of $\rho$ after $v_s$ in that set, let $f=\langle x,y\rangle$ be the
edge of $\rho$ entering $y$, and let $\omega$ be the part of $\rho$ from $v_s$ to $x$. Thus $\omega$ meets $\pi_s$ only
at $v_s$; it is empty, with $x=v_s$ and $f=g$, when $g$ joins $v_s$ to $y$.

The walk $\omega$ meets no vertex $v_k$ with $k>s$. Indeed, $\omega$ followed by $f$ starts with $g$, and if it met
such a vertex, then $\pi_s$, followed by $\rho$ up to its first such vertex $v_k$, followed by the tail of $\pi$ from
$v_k$, would be a ray beginning with $\pi_s$ and $g$. Hence, if $y=v_i$, the walk
\[
\pi_i,\ f,\ \omega\text{ reversed},\ v_s h_s v_{s+1}h_{s+1}\cdots
\]
is a ray: $\pi_i$ and the tail of $\pi$ from $v_s$ are disjoint, $\omega$ meets $\pi$ only at $v_s$, and $x$ lies on
$\pi$ only if $\omega$ is empty. So $f$ is extendable at $i<s$. It is not $h_i$: if $\omega$ is nonempty then $x$ is
off $\pi$, and if $\omega$ is empty then $f=g$ is not an edge of $\pi$. Map $g$ to $f$.

The map is injective. Distinct edges $g$ at the same $v_s$ give distinct $f$, as the rays are edge-disjoint. Suppose
$g$ at $v_s$ and $g'$ at $v_t$, with $s<t$, give the same edge $f$, with walks $\omega$, $\omega'$. If $\omega$ is
empty, both ends of $f$ lie in $\{v_0,\dots,v_s\}$; but $\omega'$ ends at a vertex $x'$ of $f$ that is not in
$\{v_0,\dots,v_{t-1}\}$, which is impossible. If $\omega'$ is empty, both ends of $f$ lie on $\pi$, while $\omega$
nonempty has its end $x$ off $\pi$. So both are nonempty, $f=\langle x,y\rangle$ for both, and $\omega$, $\omega'$
both end at $x$. Let $z$ be the first vertex of $\omega$ that lies on $\omega'$. Then $\omega$ up to $z$, followed by
$\omega'$ backwards from $z$ to $v_t$, followed by the tail of $\pi$ from $v_t$, is a ray that begins with $\pi_s$ and
$g$: the first piece meets $\pi$ only at $v_s$, the second meets $\pi$ only at $v_t$, the two meet only at $z$, and
neither meets the tail of $\pi$ beyond $v_t$. This contradicts the choice of $g$.

Therefore $R\le B$, and \eqref{eq:BR} gives $2B\ge N(d-2)+1$. Since $B$ is an integer and $N(d-2)=2n(d-2)$ is even,
$B\ge n(d-2)+1$. Finally $\sum_{s=0}^{N-1}b_s=B-1\ge n(d-2)$.
\end{proof}

The \emph{tree of extendable walks} from $v_0$ has as nodes the finite SAWs from $v_0$ that are initial segments of
rays, each joined to its one-step extensions. A node $\pi_s$ with $s\ge1$ has exactly $b_s+1$ children, and the root
has $d$ children.

\section{Atoms}\label{sec:atoms}

Let $G$ be an infinite, connected, $\Delta$-regular, vertex-transitive, loopless multigraph, and suppose
$\lambda=\lambda(G)<\Delta$. A \emph{fragment} is a finite nonempty $S\subset V$ with $|\pd S|=\lambda$, and an
\emph{atom} is a fragment of least cardinality $a$. Atoms of this kind are the edge version of those of Watkins \cite{Wat70}. Mader \cite{Mad71} proved
that a finite connected vertex-transitive simple graph has edge connectivity equal to its degree. We prove the properties
we need, since $G$ is infinite and has multiple edges.
For an atom $A$, write $G[A]$ for the multigraph on $A$ formed by the edges of $G$ with both ends in $A$, and
$E_A(S,A\setminus S)$ for its set of edges between $S$ and $A\setminus S$.

\begin{lemma}\label{lem:atoms}
\begin{enumerate}
\item $a\ge2$, and every nonempty proper subset $S$ of an atom has $|\pd S|>\lambda$.
\item The atoms partition $V$.
\item There is $e\ge1$ such that every vertex has exactly $e$ edges leaving its atom, and $\lambda=ae<\Delta$.
\item For every atom $A$, the multigraph $G[A]$ is loopless, connected and $(\Delta-e)$-regular, and every nonempty
proper $S\subset A$ satisfies $e|S|+|E_A(S,A\setminus S)|>ae$.
\end{enumerate}
\end{lemma}

\begin{proof}
(1) A single vertex has $|\pd\{v\}|=\Delta>\lambda$, and a nonempty proper subset of an atom is too small to be a
fragment. (2) Let $A\ne B$ be atoms with $A\cap B\ne\emptyset$. Then $A\cap B$ is a nonempty proper subset of $A$, so
$|\pd(A\cap B)|>\lambda$ by (1), and submodularity, $|\pd(A\cap B)|+|\pd(A\cup B)|\le|\pd A|+|\pd B|=2\lambda$, gives
$|\pd(A\cup B)|<\lambda$, which is impossible. Automorphisms map atoms to atoms and an atom exists, so by
transitivity every vertex lies in an atom, which is unique. (3) If $x,y\in A$ and $\gamma$ is an automorphism with
$\gamma x=y$, then $\gamma A$ is an atom through $y$, so $\gamma A=A$ by (2). Hence the stabiliser of $A$ is
transitive on $A$, and the number of edges from a vertex of $A$ to $V\setminus A$ is constant on $A$, say $e$. Then
$\lambda=|\pd A|=ae$, and $e\ge1$ as $G$ is connected and infinite. (4) The degree statement follows from (3). For
nonempty proper $S\subset A$ we have $|\pd S|=e|S|+|E_A(S,A\setminus S)|$, and (1) gives the inequality. Since
$e|S|<ae$, it forces $E_A(S,A\setminus S)\ne\emptyset$, so $G[A]$ is connected.
\end{proof}

Let $H$ be the multigraph whose vertices are the atoms of $G$, with one edge for each edge of $G$ joining two
distinct atoms. It is loopless, connected, infinite and $ae$-regular. For a finite nonempty set $\mathcal S$ of atoms,
the edges of $H$ leaving $\mathcal S$ are the edges of $G$ leaving the union of $\mathcal S$, and there are at least
$\lambda=ae$ of them. So $H$ satisfies the hypothesis of Section~\ref{sec:branch} with $d=ae$.

\section{Lifting walks from the quotient}\label{sec:lift}

For an atom $A$ and $u,w\in A$ let
\[
P_{uw}=\sum_\gamma t^{|\gamma|},
\]
the sum over SAWs $\gamma$ from $u$ to $w$ in $G[A]$, where $|\gamma|$ is the number of steps; in particular
$P_{uu}=1$. When $\lambda=\Delta$ we take every atom to be a single vertex, so that $H=G$, $a=1$ and $e=\Delta$.

Fix a vertex $r$ in the atom $A_0$. Let $\varpi=A_0h_0A_1h_1\cdots A_N$ be a SAW in $H$. Each $h_s$ is an edge of $G$
from a vertex $u'_s\in A_s$ to a vertex $u_{s+1}\in A_{s+1}$; put $u_0=r$. A \emph{lift} of $\varpi$ is a choice,
for each $s<N$, of a SAW $\gamma_s$ in $G[A_s]$ from $u_s$ to $u'_s$. The walk
$\gamma_0h_0\gamma_1h_1\cdots\gamma_{N-1}h_{N-1}$ is a SAW of $G$ from $r$, because the atoms $A_s$ are distinct and
each $\gamma_s$ is self-avoiding. Distinct pairs of a walk and a lift, for all $N$ together, give distinct SAWs of
$G$: the edges of the $G$-walk that join different atoms recover $\varpi$, and the segments between them recover the
$\gamma_s$. Summing over lifts and keeping only the walks $\varpi$ in the tree of extendable walks of $H$,
\begin{equation}\label{eq:lift}
\sum_kc_kt^k\ \ge\ \sum_{N\ge1}Z_N,\qquad Z_N=\sum_\varpi\ \prod_{s<N}tP_{u_su'_s},
\end{equation}
the sum over extendable $N$-step SAWs $\varpi$ of $H$ from $A_0$.

Thus in the tree of extendable walks of $H$, the edge from a node ending at $A_s$, entered at $u_s$, to the child
through $h_s$ carries the weight $tP_{u_su'_s}$. Of the $ae-1$ edges of $H$ at $A_s$ other than the entry edge, $e-1$
start at $u_s$ and have weight $t$, and for each $w\in A_s\setminus\{u_s\}$, $e$ start at $w$ and have weight
$tP_{u_sw}$.

Write $d=ae$ and $c=(d-2)/2$. For $0\le b\le d-2$ let $W(b)$ be the least total weight of $b+1$ of the $d-1$
non-entry edges at an atom, minimised over the entry vertex and the choice of edges; it is the sum of the $b+1$
smallest weights for the worst entry vertex. All weights are positive, so $W$ is increasing. Let $\varphi=\log W$ on
$\{0,1,\dots,d-2\}$ and let $\psi$ be its lower convex envelope, the largest convex function on $[0,d-2]$ lying below
$\varphi$ at the integers.

\begin{lemma}\label{lem:weighted}
If $\psi(c)\ge0$, then there is $\kappa>0$ with $Z_{2n}\ge\kappa$ for all $n\ge1$. Consequently
$\mu(G)\ge\sqrt{\Delta-1}$.
\end{lemma}

\begin{proof}
For a node $y$ of the tree at depth less than $2n$, let $Z(y)$ be the sum, over the descendants of $y$ at depth $2n$,
of the product of the edge weights on the path from $y$ to the descendant; so $Z_{2n}$ is the value at the root. If
$y$ is not the root and ends at $v_s$, it has $b_s+1$ children $z$, joined to it by non-entry edges, and
\[
Z(y)=\sum_zw(y,z)Z(z)\ \ge\ W(b_s)\min_zZ(z).
\]
At the root, $Z\ge w_{\min}\min_zZ(z)$, where $w_{\min}>0$ is the least weight of an edge of $H$ at $A_0$ from $r$. By
induction,
\[
Z_{2n}\ \ge\ w_{\min}\ \min_\varpi\ \prod_{s=1}^{2n-1}W(b_s),
\]
the minimum over extendable $2n$-step SAWs $\varpi$ of $H$ from $A_0$. Each such walk is an initial segment of a ray,
and the numbers $b_s$ depend only on the walk, so Lemma~\ref{lem:branch} and $b_0=d-2=2c$ give
$\sum_{s=1}^{2n-1}b_s\ge2nc-2c$. If $d=2$, every nonroot node has exactly one child, $\psi(0)=\log W(0)\ge0$, and
hence $Z_{2n}\ge w_{\min}W(0)^{2n-1}\ge w_{\min}$; so we may assume $d>2$. The function $\psi$ is convex and piecewise linear. It is nondecreasing, since the
constant $\varphi(0)$ lies below $\varphi$, so that $\psi\ge\varphi(0)=\psi(0)$. Let $L$ be its largest slope. With
$\bar b$ the average of $b_1,\dots,b_{2n-1}$, we have $\bar b\ge c-c/(2n-1)$, and by Jensen's inequality
\[
\sum_{s=1}^{2n-1}\varphi(b_s)\ \ge\ \sum_{s=1}^{2n-1}\psi(b_s)\ \ge\ (2n-1)\psi(\bar b)\ \ge\
(2n-1)\psi(c)-Lc\ \ge\ -Lc.
\]
So $Z_{2n}\ge w_{\min}e^{-Lc}=:\kappa$. By \eqref{eq:lift}, $\sum_kc_kt^k=\infty$, and $\mu(G)\ge1/t$.
\end{proof}

When $\lambda=\Delta$, every weight is $t$, $W(b)=(b+1)t$, and $\varphi$ is concave; so $\psi$ is the chord from $0$
to $\Delta-2$ and $\psi(c)=\tfrac12\bigl(\log t+\log((\Delta-1)t)\bigr)=0$. Lemma~\ref{lem:weighted} then gives
$\mu\ge\sqrt{\Delta-1}$, which for simple graphs is \cite[Theorem~4.1(a)]{GL15}.

\section{Degree five}\label{sec:five}

\begin{proof}[Proof of Theorem~\ref{thm:main}]
By the last paragraph of Section~\ref{sec:lift} we may assume $\lambda<5$. Here $t=1/2$. By Lemma~\ref{lem:atoms},
$a\ge2$ and $ae\le4$, so $(a,e)$ is $(2,1)$, $(2,2)$, $(3,1)$ or $(4,1)$. We show $\psi(c)\ge0$ in each case, and
Lemma~\ref{lem:weighted} gives $\mu(G)\ge2$.

$(a,e)=(2,1)$. The atom is a pair joined by $4$ parallel edges, $d=2$ and $c=0$. The only non-entry edge has weight
$t\cdot4t=1$, so $\psi(0)=\log W(0)=0$.

$(a,e)=(2,2)$. The atom $\{x,x'\}$ is joined by $3$ parallel edges, $d=4$ and $c=1$. Entering at $x$, the non-entry
edges are one edge at $x$ of weight $\tfrac12$ and two edges at $x'$ of weight $t\cdot3t=\tfrac34$. So
$W(0)=\tfrac12$, $W(1)=\tfrac54$, $W(2)=2$, and
\[
\psi(1)=\min\Bigl(\log\tfrac54,\ \tfrac12\bigl(\log\tfrac12+\log2\bigr)\Bigr)=0.
\]

$(a,e)=(3,1)$. $G[A]$ is a loopless $4$-regular multigraph on three vertices, so each pair of vertices is joined by
$2$ edges, and $d=3$, $c=\tfrac12$. For $u\ne w$, $P_{uw}=2t+4t^2=2$, so every weight is $tP_{uw}=1$. Hence
$W(0)=1$, $W(1)=2$ and $\psi(\tfrac12)=\tfrac12\log2>0$.

$(a,e)=(4,1)$. $G[A]$ is a loopless $4$-regular multigraph on four vertices, and $d=4$, $c=1$. By
Lemma~\ref{lem:atoms}(4) applied to $S=\{u,w\}$, whose edges to the other two vertices number $8-2m_{uw}$ where
$m_{uw}$ is the multiplicity of the pair, we get $2+8-2m_{uw}>4$, so $m_{uw}\le2$. The four degree equations force opposite pairs to have equal
multiplicities; writing the three values as $x,y,z$, we have $x+y+z=4$ and $0\le x,y,z\le2$, so $\{x,y,z\}$ is
$\{2,2,0\}$ or $\{2,1,1\}$, each in three labelled forms. Thus $G[A]$ is either the $4$-cycle with every edge doubled or
$K_4$ with the edges of one perfect matching doubled. In the doubled $4$-cycle, $P_{uw}=8t^2=2$ for opposite $u,w$ and
$P_{uw}=2t+8t^3=2$ for adjacent ones. In $K_4$ with the matching $\{uv,wz\}$ doubled,
$P_{uv}=2t+2t^2+4t^3=2$ and $P_{uw}=t+4t^2+5t^3=\tfrac{17}8$, and likewise for the other pairs. In both cases
every weight $tP_{uw}$ is at least $1$, so $W(b)\ge b+1$ and $\psi(1)\ge\tfrac12\log3>0$.
\end{proof}

The program \texttt{delta5\_check.py} in the repository \cite{Repo} enumerates the labelled loopless multigraphs of
each of the four types that satisfy Lemma~\ref{lem:atoms}(4), computes the path sums exactly by listing the SAWs, and
returns the least value of $\psi(c)$ over the shapes and entry vertices: $0$, $0$, $\tfrac12\log2$ and
$\tfrac12\log3$ for the four types, with one, one, one and six labelled shapes.

\section{Higher degree}\label{sec:remarks}

Sections~\ref{sec:branch}--\ref{sec:lift} hold for every $\Delta$: $\mu(G)\ge\sqrt{\Delta-1}$ whenever every atom type
$(a,e,G[A])$ allowed by Lemma~\ref{lem:atoms} has $\psi(c)\ge0$. The criterion is not always met. For atoms of size
two with $e=2$ joined by $m=\Delta-2$ edges, the weights are $t$ once and $mt^2$ twice, and
$\psi(1)=\tfrac12\log\bigl(t(t+2mt^2)\bigr)$. With $s=\sqrt{\Delta-1}$ we have
$t(t+2mt^2)=(2s^2+s-2)/s^3$, and this is less than $1$ exactly when $(s-2)(s^2-1)>0$, that is, when $\Delta\ge6$. In
that case a tree in which nodes with one child alternate with nodes having all three children satisfies
Lemma~\ref{lem:branch} and has weighted growth below one, so for $\Delta\ge6$ the branching inequality in the quotient
does not suffice on its own, and the question of Grimmett and Li remains open.

Restricting self-avoiding walks so that they cross each of a family of finite pieces of the graph at most once, as the
lifted walks of Section~\ref{sec:lift} cross each atom once, also appears in the work of Grimmett and Li on the Fisher
transformation \cite{GL20} and of Grimmett on Grigorchuk graphs \cite[Remark~3.4]{Gri26}. In
\cite[Proposition~7.4]{GL20}, and in Grimmett's bound for line graphs \cite[Theorem~7.2]{Gri26b}, walks of a base graph
are lifted to the Fisher graph or the line graph with weights given by path counts inside the replaced pieces.

\section*{Acknowledgements and funding}
This work was done by AI agents, Claude (Anthropic) and Codex (OpenAI), connected to Math Brain~\cite{MathBrainRepo}
and working together inside the Brain framework~\cite{BrainFramework}, a shared, persistent memory in which the agents
record claims, proofs, failed attempts and methods, retrieve one another's results, coordinate who works on what, and
audit each other's work. For more on the Brain and Math Brain, see the Math Brain paper~\cite{MathBrainPaper}, the Math
Brain repository~\cite{MathBrainRepo} and the Brain framework repository~\cite{BrainFramework}.

The research presented in this paper was supported by the European Research Council (ERC) under the European Union's
Horizon Europe research and innovation programme (grant agreement No.~101041711), by the Simons Foundation as part of
the Collaboration on the Mathematical and Scientific Foundations of Deep Learning, by Heights Labs, by Convex Nexus
Capital, by the Israel Science Foundation (grant number 2258/19), and by the Israel Science Foundation (ISF Grant
4101/25). Computational resources were provided in part by Google Cloud through its education credits programme
(credit 530550103).

\section*{Declaration of competing interest}
The author holds equity in Heights Labs and in Convex Nexus Capital and has a management role at Convex Nexus Capital.
These entities are included in the funding declaration above.

\section*{Data availability}
The program described at the end of Section~\ref{sec:five} is in the repository~\cite{Repo}.

\section*{Declaration of generative AI and AI-assisted technologies}
This paper is part of a broader research project on AI proving problems. Claude agents (Anthropic) proved
Theorem~\ref{thm:main}, wrote the program of Section~\ref{sec:five} and prepared the manuscript; a Codex agent
(OpenAI) checked the proof and supplied the direct form of the proof of Lemma~\ref{lem:branch}. The branching
inequality, the bridge graph and the question are due to Grimmett and Li~\cite{GL15,GL19}, lifts weighted by path counts inside pieces to Grimmett and
Li and to Grimmett~\cite{GL20,Gri26b}, atoms, in their vertex form, to Watkins~\cite{Wat70}, and the equality of edge connectivity and degree for
finite vertex-transitive simple graphs to Mader~\cite{Mad71}. The use of
edge-connectivity atoms together with the weighted branching criterion of Lemma~\ref{lem:weighted} is the
contribution here. The author directed the work and is responsible for the content of the paper
and for decisions concerning its dissemination.

\begin{thebibliography}{99}
\bibitem{GL13} G.~R. Grimmett and Z.~Li, \emph{Counting self-avoiding walks}, preprint, 2013, \arxiv{1304.7216}.
\bibitem{GL15} G.~R. Grimmett and Z.~Li, \emph{Bounds on connective constants of regular graphs}, Combinatorica
\textbf{35} (2015), 279--294, \doi{10.1007/s00493-014-3044-0}; \arxiv{1210.6277}.
\bibitem{GL19} G.~R. Grimmett and Z.~Li, \emph{Self-avoiding walks and connective constants}, in \emph{Sojourns in
Probability Theory and Statistical Physics III}, Springer Proc. Math. Stat. \textbf{300}, Springer, Singapore, 2019,
215--241, \doi{10.1007/978-981-15-0302-3_8}; \arxiv{1704.05884}.
\bibitem{GL20} G.~R. Grimmett and Z.~Li, \emph{Cubic graphs and the golden mean}, Discrete Math. \textbf{343} (2020),
111638, \doi{10.1016/j.disc.2019.111638}; \arxiv{1610.00107}.
\bibitem{Gri26} G.~R. Grimmett, \emph{Connective constants of Grigorchuk graphs}, preprint, 2026, \arxiv{2608.19349}.
\bibitem{Gri26b} G.~R. Grimmett, \emph{SAWs, trails, kites, patterns}, preprint, 2026, \arxiv{2611.00816}.
\bibitem{HM54} J.~M. Hammersley and K.~W. Morton, \emph{Poor man's Monte Carlo}, J. Roy. Statist. Soc. Ser. B
\textbf{16} (1954), 23--38, \doi{10.1111/j.2517-6161.1954.tb00145.x}.
\bibitem{Mad71} W.~Mader, \emph{Minimale $n$-fach kantenzusammenh\"angende Graphen}, Math. Ann. \textbf{191} (1971),
21--28, \doi{10.1007/BF01433466}.

\bibitem{Wat70} M.~E. Watkins, \emph{Connectivity of transitive graphs}, J. Combin. Theory \textbf{8} (1970), 23--29,
\doi{10.1016/S0021-9800(70)80005-9}.
\bibitem{Repo} Y.~Shmalo, \emph{The connective constant of non-simple transitive graphs of degree five: paper and
check program}, 2026, release v1.0,
\url{https://github.com/yspennstate/connective-constant-degree-five/tree/v1.0}.
\bibitem{MathBrainRepo} Y.~Shmalo, \emph{Math Brain: public snapshot, code and tests}, 2026,
\url{https://github.com/yspennstate/math-brain-public/tree/8c9a6bd478c4ec4082347151b7ae16db50c917dd}.
\bibitem{MathBrainPaper} Y.~Shmalo, \emph{Math Brain: shared, auditable memory for multi-agent mathematical
research}, 2026, \doi{10.5281/zenodo.24150837},
\url{https://github.com/yspennstate/math-brain-public/tree/main/papers/math-brain-2026-10-24}.
\bibitem{BrainFramework} Y.~Shmalo, \emph{Brain Framework}, public software repository, 2026,
\url{https://github.com/yspennstate/brain-framework-public/tree/5629e644cdd510173ae97e9aa92642dabd22dc57}.
\end{thebibliography}

\end{document}
